% §12.3 Preservation of Relational Invariants
%
% Monograph register: structural invariants survive approximation only when every
% constrained trial cooperates, so the preservation probability is a product over
% the constrained pairs/triples and decays as the universe grows. Reflexive
% tprate^|A|; symmetric closure as the cheap repair (FPR 1-(1-eps)^2, TPR
% 1-(1-tau)^2); functional (1-eps)^{|Y|-1} per key; transitive hardest, a
% first-order heuristic ~eps per manufactured conclusion, no cheap closure. Closes
% with the model-order reading tying back to the projection fan-in of 12.2.
% Source: towell-bernoulli-relations sec:invariants. Cross-refs ax:element-indep,
% sec:relation-composition, thm:join-correlation. DEFINED: sec:relation-invariants,
% thm:rel-reflexive, thm:rel-symmetric-closure, thm:rel-functional.

\section{Preservation of Relational Invariants}
\label{sec:relation-invariants}

A relation often carries a structural invariant: reflexivity, symmetry,
single-valuedness, transitivity. An exact relation either has the property or it
does not. An approximate relation tests each pair through its own Bernoulli trial
(\Cref{ax:element-indep}), and a property that constrains \emph{many} pairs at
once survives only when all of the relevant trials cooperate. The lesson of this
section is one of fragility: the probability that an approximate relation inherits
a global invariant is a product over the constrained pairs, and it tends to zero
as the ground universe grows. Where a cheap repair exists, as for symmetry, we
exhibit it and pay its rate cost; where none does, as for transitivity, we state
the honest limiting behavior rather than force a clean formula.

Throughout, $R \subseteq A \times A$ is an exact relation with the named property
and $\tilde R$ its approximation, with false-positive rate $\fprate$,
true-positive rate $\tprate = 1 - \fnrate$, and the per-pair trials independent.

\subsection*{Reflexivity}

A relation is \emph{reflexive} when it contains the whole diagonal: $(a, a) \in R$
for every $a \in A$. The approximation is reflexive only if each of the $|A|$
diagonal pairs survives as a true positive.

\begin{theorem}[Reflexivity is preserved with vanishing probability]
\label{thm:rel-reflexive}
If $R$ is reflexive, then $\tilde R$ is reflexive with probability
$\tprate^{|A|}$, which tends to $0$ as $|A| \to \infty$ for any $\tprate < 1$.
\end{theorem}

\begin{proof}
Each diagonal pair $(a, a)$ is a true positive with probability $\tprate$, and the
$|A|$ diagonal trials are independent (\Cref{ax:element-indep}). Reflexivity holds
exactly when all of them fire, with probability $\tprate^{|A|}$, which decays
geometrically in $|A|$ since $\tprate < 1$.
\end{proof}

\subsection*{Symmetry and its closure}

A relation is \emph{symmetric} when $(a, b) \in R$ implies $(b, a) \in R$. The two
ordered trials of an unordered pair are tested independently and need not agree,
so symmetry, like reflexivity, is preserved only with a probability that vanishes
as the number of pairs grows. Unlike reflexivity, symmetry admits a cheap repair:
rather than hope the two ordered trials agree, force agreement by taking the
\emph{symmetric closure}, which admits a pair when \emph{either} ordered trial
fires.

\begin{theorem}[Symmetric closure]
\label{thm:rel-symmetric-closure}
Let $R$ be symmetric and $\tilde R$ its approximation with rates $(\fprate,
\tprate)$. The symmetric closure $\tilde R^{s}$, which admits $(x, y)$ when $(x,
y) \in \tilde R$ or $(y, x) \in \tilde R$, is a symmetric approximate relation
with false-positive rate $1 - (1 - \fprate)^2$ and true-positive rate $1 - (1 -
\tprate)^2$.
\end{theorem}

\begin{proof}
Because $R$ is symmetric, each unordered pair is tested by two independent trials,
on $(x, y)$ and on $(y, x)$ (\Cref{ax:element-indep}). The closure admits the pair
unless both trials miss. For a negative pair each trial false-positives with
probability $\fprate$, so the closure false-positives with probability $1 - (1 -
\fprate)^2$; for a positive pair each trial is a true positive with probability
$\tprate$, so the closure accepts with probability $1 - (1 - \tprate)^2$. The
closure restores the invariant deterministically, and the same doubling that
raises the false-positive rate raises the true-positive rate: at $\tprate = 0.9$
the closure true-positive rate is $1 - 0.1^2 = 0.99$.
\end{proof}

The closure is the template for repairing a fragile invariant: replace a
conjunction of hopeful trials by an operator that enforces the property outright,
and read off the rate it costs.

\subsection*{Single-valuedness}

A relation $f \subseteq X \times Y$ is \emph{functional} (single-valued, a partial
function) when each key carries at most one value. Approximation breaks this in
only one direction. A false negative on a key's true value leaves that key with no
value, which is still single-valued; only a \emph{false positive}, which adds a
spurious second value, violates functionality. The property survives exactly when
no key acquires a spurious value.

\begin{theorem}[Single-valuedness is preserved with vanishing probability]
\label{thm:rel-functional}
Let $f \subseteq X \times Y$ be functional. A key with a defined value stays
single-valued with probability $(1 - \fprate)^{|Y| - 1}$, since each of the $|Y| -
1$ wrong values is an independent false positive at rate $\fprate$. Over the $d$
keys with a defined value, the probability that $\tilde f$ stays single-valued is
$(1 - \fprate)^{d\,(|Y| - 1)}$, which tends to $0$ as the value-universe or the
domain grows.
\end{theorem}

\begin{proof}
Fix a key $x$ with true value $y_x$. The $|Y| - 1$ pairs $(x, y)$ with $y \neq
y_x$ are all negative, each reported with probability $\fprate$ independently, so
$x$ stays single-valued with probability $(1 - \fprate)^{|Y| - 1}$. Losing $y_x$,
an event of probability $\fnrate$, leaves $x$ with no value and does not violate
the property, so it does not enter the product. Distinct keys index disjoint pair
sets, so the events are independent and the joint probability is the product over
the $d$ defined keys.
\end{proof}

\subsection*{Transitivity}

Transitivity is the hardest of the four and, unlike symmetry, has no inexpensive
closure. The constraint ranges over \emph{triples}: $(x, y) \in R$ and $(y, z) \in
R$ together require $(x, z) \in R$. A single pair participates in many triples, so
the implicated trials overlap and the preservation probability does not factor
into independent per-triple terms. The honest statement is a limiting one. There
are on the order of $|A|^3$ triple constraints, far more than the $|A|$ diagonal
constraints of reflexivity or the pair constraints of symmetry, and the trials are
coupled, so the preservation probability decays the fastest of the four. Its most
fragile part is the set of conclusion pairs $(x, z) \notin R$ that must be
\emph{manufactured} by a compensating false positive: as a first-order heuristic,
ignoring the cross-triple coupling, each such conclusion contributes a factor near
$\fprate \ll 1$, and the true probability is smaller still once the coupling is
counted. A transitive-closure repair is far more involved than the symmetric one:
adding a manufactured pair can create new required conclusions in a cascade with no
fixed per-pair cost, so we leave transitivity as the honest limiting result rather
than a doubled rate.

\subsection*{The model-order reading}

The four results share a reading in the language of \Cref{ax:element-indep}. Each
invariant partitions the universe into pairs or triples that play different error
roles: diagonal versus off-diagonal, positive versus negative, premise versus
conclusion. An approximation forced to honor the invariant is therefore a
higher-order Bernoulli relation whose blocks carry distinct rates, the same
model-order structure projection produced through its varying fan-in
(\Cref{sec:relation-composition}). Element-wise independence keeps the parameter
count linear even as the order grows, so the cost of an invariant is charged to
the model order, not to an explosion of free rates.
